\documentclass[12pt, a4paper]{article}
\usepackage{xurl}
\usepackage{hyperref}

\begin{document}
        
    \begin{itemize}
        \item A. Hogan et al., “Knowledge Graphs.” Available: \url{https://arxiv.org/pdf/2003.02320}
        \begin{itemize}
            \item This is an introduction and in-depth explanation of knowledge graphs, their history, how they work, and how they are formatted. 
            This is an important resource in getting the fundamental basics of knowledge graphs and other similar topics. 
        \end{itemize}
    \end{itemize}
    \begin{itemize}
        \item K. Liu, F. Wang, Z. Ding, S. Liang, Z. Yu, and Y. Zhou, “A review of knowledge graph application scenarios in cyber security.” 
        Available: \url{https://arxiv.org/ftp/arxiv/papers/2204/2204.04769.pdf}
        \begin{itemize}
            \item This paper goes over multiple different cyber security knowledge graphs and how they can be applied for different aspects of cybersecurity. 
            It contains resources and links to multiple open-source cyber security knowledge graphs, 
            which helped me explore how KGs are currently being used in the cybersecurity field. 
        \end{itemize}
    \end{itemize}
    \begin{itemize}
        \item A. Piplai, S. Mittal, A. Joshi, T. Finin, J. Holt, and R. Zak, “Creating Cybersecurity Knowledge Graphs From Malware After Action Reports,” 
        IEEE Access, vol. 8, pp. 211691-211703, 2020, doi: \url{https://doi.org/10.1109/access.2020.3039234}.
        \begin{itemize}
            \item This paper provides an example of a specific application of a cyber security knowledge graph through After Action Reports, 
            which describes how malware was used in attacks and how that affected a system. 
            A knowledge graph was developed from extracting information from these after-action reports to allow for reasoning capabilities 
            to determine similarities between different attacks. 
        \end{itemize}
    \end{itemize}
    \begin{itemize}
        \item Z. Wang, H. Zhu, P. Liu, and L. Sun, “Social engineering in cybersecurity: a domain ontology and knowledge graph application examples,” 
        Cybersecurity, vol. 4, no. 1, Aug. 2021, doi: \url{https://doi.org/10.1186/s42400-021-00094-6}.
        \begin{itemize}
            \item This paper develops an ontology and knowledge graph based on social engineering attacks. 
            I was reading this paper while trying to figure out what my specific topic was going to be.
        \end{itemize}
    \end{itemize}
    \begin{itemize}
        \item K. Liu, F. Wang, Z. Ding, S. Liang, Z. Yu, and Y. Zhou, “A review of knowledge graph application scenarios in cyber security.” 
        Available: \url{https://arxiv.org/ftp/arxiv/papers/2204/2204.04769.pdf} 
        \begin{itemize}
            \item This was another survey paper going over the many applications of knowledge graphs in cyber security. 
            It also pointed me to other uses of knowledge graphs in the cyber security field
        \end{itemize}
    \end{itemize}
    \begin{itemize}
        \item Y. Ma et al., “The Advancement of Knowledge Graphs in Cybersecurity: A Comprehensive Overview,” Mechanisms and machine science (Print), 
        pp. 65-103, Dec. 2023, doi: \url{https://doi.org/10.1007/978-3-031-42987-3\_6}. 
        \begin{itemize}
            \item This was another survey paper going over the many applications of knowledge graphs in cyber security. 
            It also pointed me to other uses of knowledge graphs in the cyber security field
        \end{itemize}
    \end{itemize}
    \begin{itemize}
        \item C. Eagle and K. Nance, The Ghidra Book: The Definitive Guide. San Francisco, CA, USA: No Starch Press, 2020.
        \begin{itemize}
            \item This book is an extensive guide on how to use Ghidra, a software reverse engineering tool that analyzes binary and executable files. 
            This was an important resource for me to start learning how Ghidra works. 
        \end{itemize}
    \end{itemize}
    \begin{itemize}
        \item J. Huang, J. Zhang, J. Liu, C. Li, and R. Dai, “UQuery: Static Security Analysis of PHP-Based Web Programs Using Graph Models,” 
        2024 IEEE Conference on Communications and Network Security (CNS), pp. 1-9, Sep. 2024, doi: \url{https://doi.org/10.1109/cns62487.2024.10735619}. 
        \begin{itemize}
            \item This paper goes over a framework that develops dependency graphs used for server-side web applications. 
            Its two main applications are detecting unrestricted file upload vulnerabilities and identifying information leakage through exploitable race conditions. 
            UQuery was able to discover and report 6 new CVEs through its dependency graphs. One of the authors of this paper helped me decide my research topic.
        \end{itemize}
    \end{itemize}
    \begin{itemize}
        \item G. Agrawal, K. Pal, Y. Deng, H. Liu, and C. Baral, “AISecKG: Knowledge Graph Dataset for Cybersecurity Education.” 
        Available: \url{https://ceur-ws.org/Vol-3433/paper6.pdf} 
        \begin{itemize}
            \item This paper goes over AISecKG, which is an AI-enabled automated KG for cybersecurity education. 
            Since there are many aspects of cybersecurity, including different attacks, the tools used by both attackers and defenders, and more, 
            developing knowledge graphs for different scenarios can help students better understand certain CS concepts better. 
            The paper explains how they trained a language model to extract named entities related to cybersecurity to create the dataset used for KG development.
        \end{itemize}
    \end{itemize}
    \begin{itemize}
        \item M. McCain, “A Small Survey on Generative Artificial Intelligence and Machine Learning Techniques in Software Reverse Engineering,” unpublished
        \begin{itemize}
            \item This survey paper shows use cases for generative AI and machine learning for software reverse engineering. 
            It was useful for getting an understanding of software reverse engineering and what kind of projects are created in that field, 
            and it cited some other useful papers for my research. 
        \end{itemize}
    \end{itemize}
    \begin{itemize}
        \item Z. Chen, B. Pan, and Y. Sun, “A Survey of Software Reverse Engineering Applications,” Lecture Notes in Computer Science, 
        pp. 235-245, 2019, doi: \url{https://doi.org/10.1007/978-3-030-24268-8\_22}. 
        \begin{itemize}
            \item The survey of software reverse engineering applications paper shows use-cases and applications for software reverse engineering, 
            including malware analysis.
        \end{itemize}
    \end{itemize}
    \begin{itemize}
        \item Vishal Karande, S. Chandra, Z. Lin, J. Caballero, L. Khan, and K. Hamlen, “BCD,” pp. 393-398, May 2018, doi: \url{https://doi.org/10.1145/3196494.3196504}. 
        \begin{itemize}
            \item This paper develops a decomposition graph based on functions to break them down into different components 
            to better understand how they work and to try to get back some of the modularity lost during the initial compilation process. 
            This was an interesting way of showing how to gain back some modularity that was lost but doesn't work if the data was obfuscated 
            (at least it was only tested when the data was unobfuscated). It was also tested with obfuscated data 
            (aka, testing using only call-graph and data-reference graphs, and its accuracy was substantially lower). 
        \end{itemize}
    \end{itemize}
    \begin{itemize}
        \item S.-Y. Yu, Achamyeleh, Yonatan Gizachew, C. Wang, A. Kocheturov, P. Eisen, and M. Abdullah, 
        “CFG2VEC: Hierarchical Graph Neural Network for Cross-Architectural Software Reverse Engineering,” arXiv.org, 2023. \url{https://arxiv.org/abs/2301.02723}
        \begin{itemize}
            \item This paper introduces a neural network graph to help improve the tools for software reverse engineering for things like predicting function names. 
            It creates something called a graph-of-graph, which combines the information from Ghidra's control flow graphs and function call graphs, 
            then the neural network uses that to learn the representation for the functions in the binary for function name predictions. 
            This paper helps show how to make another kind of graph (not knowledge graph) based on Ghidra analysis information, 
            specifically the control flow graph and function call graphs.
        \end{itemize}
    \end{itemize}
    \begin{itemize}
        \item C. Shimizu, R. McGranaghan, A. Eberhart, and A. C. Kellerman, “Towards a Modular Ontology for Space Weather Research,” arXiv.org, 2020. 
        \url{https://arxiv.org/abs/2009.12285} (accessed Jan. 25, 2026). 
        \begin{itemize}
            \item This paper proposes a modular ontology for space weather research. While the content of the research is not related to my own, 
            this paper explains how to structure an ontology (or at least how they structured theirs), 
            which helped me figure out how to structure my own schema and ontology.
        \end{itemize}
    \end{itemize}
    \begin{itemize}
        \item A. Ellis, S. Severyn, F. Novakazi, H. Banaee, and C. Shimizu, “OntoPret: An Ontology for the Interpretation of Human Behavior,” 
        arXiv.org, 2025. \url{https://arxiv.org/abs/2510.23553} (accessed Jan. 25, 2026). 
        \begin{itemize}
            \item This paper analyzes and interprets human behavior, specifically looking for deviations or deception, 
            to better help machines engage with humans in frameworks in the 5.0 industry. 
        \end{itemize}
    \end{itemize}
    \begin{itemize}
        \item S. N. Severyn, “Adapting Linguistic Deception Cues for Malware Detection,” Ohiolink.edu, 2025. 
        \url{https://etd.ohiolink.edu/acprod/odb\_etd/etd/r/1501/10?clear=10&p10\_accession\_num=wright1421025881} (accessed Jan. 25, 2026). 
        \begin{itemize}
            \item This paper detects patterns similar to human deception in disassembled malware to help determine if a piece of software was malware or benign. 
            It compares the patterns of deception within human language with the patterns of behavior that exist in malware. 
            The ways that malware behavior was defined and classified in this thesis helped me find what patterns and behaviors to look for to detect malware
            in my own research. 
        \end{itemize}
    \end{itemize}
    \begin{itemize}
        \item A. Eberhart, C. Shimizu, S. Chowdhury, Md. K. Sarker, and P. Hitzler, “Expressibility of OWL Axioms with Patterns,” 
        Lecture Notes in Computer Science, pp. 230-245, 2021, doi: \url{https://doi.org/10.1007/978-3-030-77385-4\_14}. 
        \begin{itemize}
            \item Proposes plugin for protege to express axioms from a list of basic types to make it easier to input simple axioms into an Ontology. 
            This paper lists the 17 types of axioms that I used in the axiom specification for my ontology. 
        \end{itemize}
    \end{itemize}
    \begin{itemize}
        \item “Reverse Engineering: Techniques, Applications, Challenges, Opportunities,” 
        International Research Journal of Modernization in Engineering Technology and Science, Aug. 2024, doi: \url{https://doi.org/10.56726/irjmets60817}. 
        \begin{itemize}
            \item This paper provides an overview of reverse engineering, its applications, and the relating legal issues. 
            This paper was used to identify some of the potential ethical and legal issues of my research.
        \end{itemize}
    \end{itemize}
    \begin{itemize}
        \item J. Acosta and D. Krych, “Hands-on Cybersecurity Studies: Uncovering and Decoding Malware Communications- Malware Analysis with Ghidra.” 
        Available: \url{https://apps.dtic.mil/sti/trecms/pdf/AD1119396.pdf }
        \begin{itemize}
            \item A report of an exercise that helps show how to determine how a certain kind of malware communicates across a private network 
            and analyzing that communication through Ghidra. 
        \end{itemize}
    \end{itemize}
    \begin{itemize}
        \item A. Singhal and Saathwick Venkataramalingam, “Malware Analysis and Reverse Engineering: Unraveling the Digital Threat Landscape,” 
        Social Science Research Network, Jan. 2023, doi: \url{https://doi.org/10.2139/ssrn.4649754}. 
        \begin{itemize}
            \item This paper gives a background of what malware is, the common types of malwares, how to analyze malware, and the steps of reverse engineering malware. 
            It also mentions the tools commonly used at the time for analyzing and reverse engineering malware. 
            It uses the tools and strategies mentioned to analyze a known ransomware file to determine how dangerous it is. 
            The authors of the paper obtained the ransomware from an open-source repository of malware used for academic purposes,
            which is the same method I used to obtain the executables for my research.
        \end{itemize}
    \end{itemize}
    \begin{itemize}
        \item “CS6038/CS5138 Malware Analysis, UC by ckane,” Malware.re, 2021. \url{https://class.malware.re/2021/03/08/ghidra-scripting.html} 
        \begin{itemize}
            \item A helpful resource I found in learning the basics of scripting in Ghidra. 
            It goes over demos of how to make a basic Ghidra script, and eventually how to work with the decompiled code of a file. 
            This demo does have some outdated code since it used Jython (which is stuck at version 2.7) 
            instead of the updated PyGhidra (which supports python version 3), so while the code itself might not be what I am looking for, 
            it was still a helpful resource in understanding the basics of how ghidra scripting works. 
        \end{itemize}
    \end{itemize}
    \begin{itemize}
        \item R. Tamilkodi, V. B. Sankar, Shrija Madhu, L. Revathi, V. Sai, and Ernagula Ajay, 
        “Exploring IoT Device Vulnerabilities Through Malware Analysis and Reverse Engineering,” pp. 279-283, Nov. 2024, 
        doi: \url{https://doi.org/10.1109/icdici62993.2024.10810929}. 
        \begin{itemize}
            \item This paper describes the analysis, detection, and reverse engineering techniques that can be used to detect malware within IoT devices.
            The proposed process of malware detection includes collecting the malware sample data from the IoT devices, 
            conducting both static and dynamic analysis to identify the vulnerabilities, 
            then using heuristic and signature-based detection to generate more malware signatures to put into a database to be referenced later. 
            Mitigation strategies are then developed to implement back into the IoT devices, then the cycle continues. 
        \end{itemize}        
    \end{itemize}

\end{document}

